% Nouveau EGA V (cote 25): archive cote 25, site page 31 (PDF page 32).
% typescript. Draft compared: none.
% AI-transcribed by Claude Opus 5.5 (claude-opus-5-5, adaptive thinking, effort high), Message Batches API, batch ega25b (msgbatch_01A7n5YuXubUPNK1zsepi8RG), 2026-11-13; original AI transcription, subsequently checked against both supplied scan parts; see PAGES_CHECKED.tsv and CHECK_DECISIONS.tsv.
\scanpage{25/31}
\mspage{\hand{27}}

8. \uline{Dimension de l'ensemble des hyperplans exceptionnels}.

8.1.Dans les n°s précédents, et notamment 2 et 3, nous avons donné des
énoncés affirmant que l'ensemble des \hand{$\xi\in$} $\overset{\hand{\vee}}{P}$ tels que Y ait une certaine
propriété P est constructible, et contient le point générique $\eta$,  ou
encore que l'ensemble $Z_P$ des \hand{$\xi\in$} $\overset{\hand{\vee}}{P}$ ``exceptionnels pour P'' est construc-%
tible, et est rare, i.e. son adhérence est de codimension \hand{$\geq$} \hand{1}) \placedmargin{\unsure{.}}
(NB on suppose S=Spec(k)). Dans certains cas on peut préciser cet énoncé
en donnant une meilleure minoration de cette codimension, ce qui est impor-%
tant dans diverses questions. Par exemple, lorsqu'on sait que cette codimen
sion est \hand{$\geq$} 2, il s'ensuit qu'une droite ``assez générale'' D de $\overset{\hand{\vee}}{P}$
ne rencontre pas $Z_P$ , d'où l'existence (si k est infini) de ``pinceaux
linéaires'' de sections hyperplanes \hand{\unsure{Dt}} $Y_{\hand{\xi}}$ (\hand{$\xi$} un point géométri
que de D) qui sont toutes la propriété P (cf n° $-\ -\ -\ -$ pour des exem-%
ples. - Du point de vue rédaction, comme les résultats du présent n°
précisent certains résultats des n°s antérieurs, la question se pose s'il
y a bien lieu de faire ce ratrappage dans un n° à part, ou de donner la
version précisée à mesure. Redactor decidetur.



\mspage{\hand{31}}
